\section{本章小结}
导航把“我在哪、哪里能走、怎样到达、失败怎么办”连接成闭环。TF 统一坐标，地图表示环境，代价地图加入障碍与外形约束，规划器生成路径，控制器跟踪路径并持续处理局部变化。应熟练定位断点和读参数；SLAM、控制器插件与行为树扩展按具体方向深入。

\begin{problemset}
\item 基础：系统已有全局路径，但控制器未输出速度，应先检查规划器还是控制链？
\item 基础：map、odom、base\_link 的典型作用是什么？
\item 基础：分辨率 0.1 m/cell、原点为 (1,2)、无旋转，索引 (2,3) 的格心在哪里？
\item 基础：静态层、障碍层、膨胀层分别承担什么？
\item 基础：A-B 代价 2、B-C 代价 2、A-C 代价 5，Dijkstra 选择哪条路线？
\item 基础：Fallback 的第一个子项返回 Running，是否立即执行下一个子项？
\item 提高：机器人速度 2 m/s、可用制动减速度 1 m/s²，忽略其它因素的制动距离是多少？为什么局部地图边界不能小于它？
\item 实践：修改本章最短路程序，使 A-B=4、B-C=4、A-C=5，运行并记录到 C 的距离。
\end{problemset}

\section{本章验收}
\begin{enumerate}
    \item 对照实际 Nav2 配置或离线日志，展示规划器、控制器与代价地图的输入输出，解释一个故障的排查路径。
    \item 演示独立 TF 树查询，说明 map、odom、base\_link 的关系；在模拟环境中展示一处变换缺失的表现。
    \item 按检查者给定的失败序列演示行为树或等价模拟器，展示至少两种恢复分支、重试上限和最终失败出口。
\end{enumerate}
